% ****** Start of file apssamp.tex ******
%
%   This file is part of the APS files in the REVTeX 4.2 distribution.
%   Version 4.2a of REVTeX, December 2014
%
%   Copyright (c) 2014 The American Physical Society.
%
%   See the REVTeX 4 README file for restrictions and more information.
%
% TeX'ing this file requires that you have AMS-LaTeX 2.0 installed
% as well as the rest of the prerequisites for REVTeX 4.2
%
% See the REVTeX 4 README file
% It also requires running BibTeX. The commands are as follows:
%
%  1)  latex apssamp.tex
%  2)  bibtex apssamp
%  3)  latex apssamp.tex
%  4)  latex apssamp.tex
%
\documentclass[%
 reprint,
%superscriptaddress,
%groupedaddress,
%unsortedaddress,
%runinaddress,
%frontmatterverbose, 
%preprint,
%preprintnumbers,
%nofootinbib,
%nobibnotes,
%bibnotes,
 amsmath,amssymb,
 aps,
%pra,
%prb,
%rmp,
%prstab,
%prstper,
%floatfix,
]{revtex4-2}

\usepackage{graphicx}% Include figure files
\usepackage{dcolumn}% Align table columns on decimal point
\usepackage{bm}% bold math
\usepackage[utf8]{inputenc}
\usepackage[spanish]{babel}
%\usepackage{hyperref}% add hypertext capabilities
%\usepackage[mathlines]{lineno}% Enable numbering of text and display math
%\linenumbers\relax % Commence numbering lines

%\usepackage[showframe,%Uncomment any one of the following lines to test 
%%scale=0.7, marginratio={1:1, 2:3}, ignoreall,% default settings
%%text={7in,10in},centering,
%%margin=1.5in,
%%total={6.5in,8.75in}, top=1.2in, left=0.9in, includefoot,
%%height=10in,a5paper,hmargin={3cm,0.8in},
%]{geometry}

\begin{document}

\preprint{APS/123-QED}

\title{T\'itulo de la pr\'actica}% Force line breaks with \\
%\thanks{A footnote to the article title}%

\author{Apellido1, Nombre1.}
\author{Apellido2, Nombre2.}
\author{Apellido3, Nombre3.}
 \affiliation{Universitdad Tecnol\'ogica de M\'exico, Campus Le\'on.}%Lines break automatically or can be forced with \\
 %t\email{Second.Author@myuniinstitution.edu}

\date{\today}% It is always \today, today,
             %  but any date may be explicitly specified

\begin{abstract}
Es un enunciado resumido del  contenido de  la  practica y  organizado secuencialmente  (menos de 100 palabras). Establece brevemente el problema y propósito de la práctica. Indica el plan teórico o experimental a que se siguió. Resume las principales aportaciones y principales conclusiones.\\ \textbf{Importante}:El  resumen  es  la  información  condensada  que  el  profesor  recibe  de  la  práctica realizada.  Debe  ser  corto  y  claro  para  que  el  profesor  determine  si  se  entendió  y  realizo correctamente  la  práctica.  Incluye  información  de  seguridad  cuando  sea  necesario.  No  evalúes conclusiones.\\\textbf{Súper  Importante}:Escribe  el  resumen  al  último  para  estar  seguro  que  refleja  con  exactitud  la práctica realizada
\begin{description}
\item[Palabras clave]: medición, vernier.
\end{description}
\end{abstract}

%\keywords{Suggested keywords}%Use showkeys class option if keyword
                              %display desired
\maketitle

%\tableofcontents

\section{\label{sec:introduccion}Introducción}

 La  introducción  debe  responder  a  la  pregunta  de  "porqué  se  ha  hecho  este  trabajo".Una  buena introducción,  es  una  oración  clara  del  problema  y  de  las  razones  por  las  que  lo  estamos estudiando.  Nos  da  una  concisa  y  apropiada  discusión  del  problema,  su  significado,  alcances  y limitaciones.\\\textbf{Importante}:En esta sección se responde a la pregunta de "cómo se ha hecho la práctica".\\\textbf{Súper Importante}:Es conveniente que el último párrafo de la Introducción se utilice para resumir el objetivode la práctica.

\section{Objetivo y planteamiento del problema}

Identificación del problema central de estudio y enfoque de la práctica.Importante:Esta  sección  se  relaciona  con  el  procedimiento  experimental  presentado  en  el protocolo, pero no lo es del todo. Se trata que identifiques dentro de tu desarrollo experimental el sentido que tuvo la práctica para poder abordar el tema a estudiar.

\section{Desarrollo experimental}

En esta secci\'on deben describirse los materiales y métodos utilizados para el desarrollo de la práctica. Para describir el arreglo experimental, puede utilizarse una lista:
\begin{itemize}
    \item Elemento 1
    \item Elemento 2
    \item \ldots
\end{itemize}
o puede describirse en un párrafo cómo se formó el arreglo. Sin embargo, es de carácter obligatorio que el arreglo experimental se describa en un diagrama (no fotos).
Finalmente, en esta secci\'on debe describirse cómo se llevó a cabo el experimento.
\section{Discusión y Resultados}
Conjunto  de  datos  obtenidos  experimentalmente  y  tratados  estadísticamente.  Usa  tablas  paraorganizar y resumir los resultados\\\textbf{Importante}:Incluye  solo  los  datos  importantes  y  relevantes,  pero  suficientes  para  justificar  tus conclusiones. Usa ecuaciones, gráficos y figuras.\\\textbf{Súper  Importante}:Si  cuentas  con  datos  teóricos  no  olvides  obtener  el  error  experimental  (\%E). Esta información es muy relevante para justificar tus conclusiones.
Cuando construyas tablas no olvides:
\begin{enumerate}
    \item  Titulo de la tabla
    \item Distinción clara de celdas.
    \item Títulos de columnas.
    \item Notas al pie de la tabla
    \item Numeración de tablas
\end{enumerate}
\section{Conclusiones}
Las conclusiones en un reporte deberán recoger los puntos primordiales del trabajo realizado, siendo una de las partes de un informe importante.

\section{Bibliografía}

En esta sección se anotan todas las fuentes bibliográficas consultadas para el desarrollo de la práctica.

%\bibliography{apssamp}% Produces the bibliography via BibTeX.

\end{document}
%
% ****** End of file apssamp.tex ******